%!TEX root = ../../thesis.tex

\section{Input Actions}\label{sec:targets-input-actions}

Since the main goal is evaluating applications of \gls{rl} to \gls{cgf}, the input
action model targets serve mainly as baseline checks for the correctness of our
EfficientZero learner.
We include three targets: \emph{high-and-low} poses a pure bandit-like problem to the learner,
while \emph{sequence} and \emph{open62541} require the learner to solve a full POMDP problem.

\subsection{High-and-low}\label{sec:target-high-and-low}

This target represents a bandit-like example.\index{target!high-and-low}
It reads a single byte and depending on its value, executes either a high-coverage
or low-coverage region.
Importantly, it does not change program state.
This makes it easy to learn since EfficientZero must only internalize what inputs
produce what coverage.
The reward is marginal coverage.
\Cref{lst:target-high-and-low} shows an excerpt of the program.

\begin{ccode}[float=htbp,
  caption={[The \emph{high-and-low} target]%
    An excerpt of the \emph{high-and-low} target.
    Each byte selects either a high- or low-coverage function without changing
    program state. The loop is bounded at 128 iterations.},
  label={lst:target-high-and-low}]
int main(void) {
    char c;
    int i = 0;

    while ((c = getchar()) != EOF && i++ < 128) {
        if (c >= 'A' && c <= 'D') {
            high_coverage();
        } else {
            low_coverage();
        }
    }

    return EXIT_SUCCESS;
}
\end{ccode}

As described in~\cref{sec:input-level-actions}, we model actions as actual PUT input
and a single environment step then consists of supplying a single byte to the PUT.
The action space therefore has a size of $\lvert\mathcal{A}\rvert=256$.
Observations are the marginal block hit counts.

\subsection{Sequence}\label{sec:target-sequence}

The \emph{sequence} target is an artificial 16-byte passcode check implemented in C.\index{target!sequence}
It contains a fixed passcode and reads input byte-by-byte from standard input.
Whenever a wrong byte is provided, the program terminates.
Otherwise, at the end, the program executes a high-coverage function.
The main goal is for the learner to successively learn the passcode by following
the non-terminating, coverage-increasing path of choosing correct bytes.
The reward in this case is a flat survival bonus per step plus a final coverage
reward on termination.
\Cref{lst:target-sequence} shows a program excerpt.

\begin{ccode}[float=htbp,
  caption={[The \emph{sequence} target]%
    An excerpt of the \emph{sequence} target.
    The read loop exits on the first byte that deviates from the passcode,
    while a correct passcode pays a high reward.},
  label={lst:target-sequence}]
const int PASSCODE_LENGTH = 16;
const char passcode[16] = "ABCDEFGHIJKLMNOP";

int main(void) {
    int c;
    int i = 0;

    while (i < PASSCODE_LENGTH && (c = getchar()) != EOF) {
        if (c != passcode[i]) {
            return EXIT_FAILURE;
        }
        i++;
    }

    high_coverage();

    return EXIT_SUCCESS;
}
\end{ccode}

The action, reward and observation setting is similar to that of \cref{sec:target-high-and-low}, with
two exceptions.
First, since we are now in a full POMDP setting instead of just a bandit, we include
the cumulative block hit counts in the observation to allow the learner to better identify
the state it is in.
Second, we restrict the action space to the 26 alphabetic characters \code{[A-Z]}
to reduce the branching factor.
Past actions are not included in the observations, so the learner has to base its
decision only on the actual blocks that were executed.

\subsection{open62541}\label{sec:target-open62541}

open62541 is an open-source OPC~UA implementation, a protocol stack of the stateful kind
described in~\cref{sec:fuzzing-stateful}.\index{target!open62541}
The PUT is its \code{tutorial\_server\_firststeps} example server, and the goal is to reach the
connected state: an activated session on an open secure channel.
We boot the server until it blocks accepting connections, inject a client connection and let the
server accept it, and snapshot there, so every episode starts with the server blocked on the first
read of an established stream.

The target is a protocol-level \emph{sequence} (\cref{sec:target-sequence}): an ordered prefix of
correct choices keeps the PUT alive, a wrong choice ends the episode, and the observation carries
the cumulative block hit counts next to those of the current step.
What differs is that an action is a whole protocol message instead of a byte.
The 15 actions are the message templates of the target's OPC~UA corpus, of which \code{hel},
\code{opn}, \code{create\_sess} and \code{activate\_sess} form the connect path, ten further channel
and session requests are self-loops the server answers without advancing towards a session, and
\code{clo} closes the channel.
Each template carries session-coupling fields, the secure channel identifier and the authentication
token, which the environment rewrites from the server's responses so that the agent learns which
message to send when rather than how to forge an identifier.

The reward is a flat survival bonus of one for every step the server survives plus ten for the step
that activates a session.
Coverage is observed but never paid for.
The episode ends when the server closes the channel after a rejected message, when a session is
activated, or at the twelve-step cap.
Because activating also terminates the episode, loitering on self-loops for the full budget returns
\num{12} while connecting immediately returns \num{14}, an undiscounted margin of two for the exact
four-message ordering, and the maximum of \num{22} requires holding the channel open for eleven
steps and activating on the twelfth.
